% !TEX root = ../proyect.tex

\chapter{Conclusiones}\label{cap11}

\section{Grado de cumplimiento}\label{sec:cumplimiento}

El proyecto BargAIn ha completado todos los bloques de desarrollo planificados:

\begin{itemize}
  \item \textbf{F1:} an\'alisis, requisitos y dise\~no de arquitectura.
  \item \textbf{F2:} infraestructura base de desarrollo y CI/CD.
  \item \textbf{F3:} backend core completo con todos los m\'odulos de dominio y API documentada.
  \item \textbf{F4:} frontend base y avanzado --- autenticaci\'on, listas, cat\'alogo, mapa,
    notificaciones, perfil y Places enrichment.
  \item \textbf{F5:} sistema de optimizaci\'on multicriterio (OR-Tools + OpenRouteService),
    scraping de precios (Scrapy + Playwright, 4 cadenas), OCR de tickets (Google Cloud Vision API
    + fuzzy matching), asistente LLM (Google Gemini 2.0 Flash Lite) y wiring de pantallas
    m\'oviles a endpoints reales.
  \item \textbf{F6:} portal business completo para PYMEs --- servicio de aprobaci\'on compartido,
    gesti\'on de precios y promociones, tests de integraci\'on, UAT verificada, notificaciones push
    y email, sincronizaci\'on documental.
  \item \textbf{F7:} pruebas integrales (unitarias, de integraci\'on y E2E), despliegue en
    staging y cierre documental de la memoria.
\end{itemize}

Adicionalmente, tras el cierre del plan inicial se incorpor\'o una fase de adaptaci\'on de la app
de usuario al uso web de escritorio (dise\~no responsive, atajos de teclado, exportaciones y
deep-linking), documentada en los Cap\'itulos~\ref{cap08} y~\ref{cap09}.

El sistema est\'a desplegado en staging (Render.com) y los cuatro flujos E2E cr\'iticos han sido
verificados de forma automatizada con Playwright.

\section{Aportaciones principales del trabajo}\label{sec:aportaciones}

\begin{itemize}
  \item \textbf{Algoritmo de optimizaci\'on multicriterio propio:} funci\'on de scoring ponderada
    (precio, distancia, tiempo) con OR-Tools para el VRP y OpenRouteService para distancias reales
    de conducci\'on. Es el n\'ucleo diferencial del TFG.

  \item \textbf{Pipeline de datos de precios automatizado:} cuatro spiders Scrapy para Mercadona,
    Carrefour, Lidl y DIA, con programaci\'on peri\'odica v\'ia Celery Beat y normalizaci\'on de
    producto contra cat\'alogo interno.

  \item \textbf{OCR sobre tickets con Google Vision API:} extracci\'on de texto, fuzzy matching
    contra cat\'alogo y conversi\'on autom\'atica a \'items de lista. Elimina la entrada manual.

  \item \textbf{Portal business para PYMEs:} flujo completo de registro, aprobaci\'on por
    administrador, gesti\'on de precios y promociones, con visibilidad para el consumidor final.

  \item \textbf{Asistente LLM con guardrails de dominio:} Gemini 2.0 Flash Lite con restricci\'on
    tem\'atica a consultas de compra, historial de conversaci\'on y rate limiting.

  \item \textbf{Arquitectura full-stack coherente:} separaci\'on clara entre apps Django, workers
    Celery independientes, PostGIS para geoespacial y React Native + Expo para iOS/Android desde
    una \'unica base de c\'odigo.

  \item \textbf{Verificaci\'on automatizada:} suite de tests backend (unitarios + integraci\'on,
    \textgreater80\% cobertura), 111 tests de frontend con Jest, tests E2E con Playwright
    (4 flujos cr\'iticos) y UAT manual en iPhone.
\end{itemize}

\section{Limitaciones}\label{sec:limitaciones}

\begin{itemize}
  \item \textbf{Cobertura de scraping:} el sistema cubre actualmente cuatro cadenas de
    supermercados (Mercadona, Carrefour, Lidl, DIA). Las cadenas con protecci\'on agresiva contra
    scrapers (Alcampo, Aldi) requieren mantenimiento adicional de los spiders.

  \item \textbf{Caducidad de precios:} los precios de scraping tienen TTL de 48 horas y los de
    crowdsourcing 24 horas; pueden no reflejar cambios de precio intradiarios.

  \item \textbf{Cold starts en Render free tier:} el plan gratuito de Render puede tener tiempos
    de arranque de hasta 30 segundos tras inactividad prolongada, lo que afecta a la primera
    petici\'on tras un per\'iodo de reposo.

  \item \textbf{Certificado iOS gratuito:} la distribuci\'on con Sideloadly y Apple ID gratuito
    genera certificados con caducidad de 7 d\'ias, que requieren re-sideload peri\'odico para el
    dispositivo de demo.

  \item \textbf{RNF-005 sin load test real:} la justificaci\'on de escalabilidad a 10.000 usuarios
    concurrentes es arquitectural (Celery workers independientes, PostGIS spatial index, API
    stateless). Un load test completo requeriría infraestructura de pago fuera del alcance del TFG.
\end{itemize}

\section{Lecciones aprendidas}\label{sec:lecciones}

\begin{itemize}
  \item \textbf{Modelo h\'ibrido backend Docker / frontend nativo en host:} determinante para la
    productividad en Windows. Docker volumes en Windows rompen el HMR de Metro/Expo; ejecutar el
    frontend nativo elimina este problema manteniendo el backend aislado.

  \item \textbf{Contratos API expl\'icitos y documentaci\'on viva:} mantener TASKS.md, ADRs y la
    memoria actualizados durante el desarrollo (no al final) reduce la fricci\'on entre backend y
    frontend y facilita la incorporaci\'on de agentes IA al flujo.

  \item \textbf{OR-Tools frente a algoritmo propio:} usar la librer\'ia de optimizaci\'on
    combinatoria de Google (producci\'on en empresas log\'isticas) es preferible a reinventar el
    VRP desde cero. El valor del TFG est\'a en la integraci\'on y la funci\'on de scoring, no en
    reimplementar un solucionador gen\'erico.

  \item \textbf{OCR cloud vs local:} Google Vision API supera a Tesseract en precisi\'on sobre
    tickets de supermercado reales (texto impreso variable), justificando el coste de la llamada
    API frente a la alternativa local.

  \item \textbf{Playwright request fixture para flujos mobile-only:} los flujos de lista y
    optimizador son exclusivos de la app m\'ovil. Cubrirlos con tests API-level (Playwright
    \texttt{request} fixture) permite automatizaci\'on completa sin necesitar una interfaz web
    consumer.
\end{itemize}

\section{Trabajo futuro}\label{sec:trabajofuturo}

\begin{enumerate}
  \item \textbf{Publicaci\'on en App Store y Google Play:} el proyecto est\'a t\'ecnicamente listo
    para iniciar el proceso de publicaci\'on. Requiere cuenta de desarrollador Apple/Google y
    revisi\'on de las gu\'ias de contenido.

  \item \textbf{Ampliaci\'on de scrapers:} incorporar Alcampo, Aldi, Lidl.es (versi\'on
    alternativa), El Corte Ingl\'es alimentaci\'on y PYMEs locales v\'ia integraci\'on API propia.

  \item \textbf{Precios en tiempo real:} varios supermercados est\'an desplegando APIs semi-p\'ublicas
    o feeds RSS de precios; la arquitectura est\'a preparada para incorporarlos como fuente
    \texttt{api} en el modelo \texttt{Price}.

  \item \textbf{Sistema de rese\~nas de productos:} los usuarios podr\'ian valorar la
    calidad/frescura de productos por tienda, a\~nadiendo una dimensi\'on cualitativa al
    comparador.

  \item \textbf{Suscripci\'on premium para PYMEs:} el modelo de negocio contempla tres tiers
    (\texttt{free}, \texttt{basic}, \texttt{premium}) con funcionalidades diferenciadas; la
    implementaci\'on actual es el tier gratuito.

  \item \textbf{Internacionalizaci\'on:} la arquitectura soporta m\'ultiples pa\'ises; la
    expansi\'on requiere scrapers por mercado y ajuste de la l\'ogica de normalizaci\'on de
    cadenas.
\end{enumerate}

\section{Cierre}\label{sec:cierre}

BargAIn demuestra que es posible construir, en el marco de un TFG, un sistema full-stack de
complejidad real que integra optimizaci\'on combinatoria, geoespacial, procesamiento de im\'agenes,
generaci\'on de lenguaje natural y scraping automatizado bajo una arquitectura coherente y
verificada. El proyecto cumple todos los objetivos planteados en la fase de an\'alisis y deja una
base t\'ecnica s\'olida para su evoluci\'on como producto real.
